\documentclass[a4paper,12pt]{article}

\usepackage[T2A]{fontenc}
\usepackage[utf8]{inputenc}
\usepackage[russian]{babel}

\usepackage{amsmath, amssymb, amsthm, mathtools}
\usepackage{helvet}
\renewcommand{\familydefault}{\sfdefault}
\usepackage{icomma}
\usepackage[left=18mm,right=18mm,top=18mm,bottom=20mm]{geometry}
\usepackage{microtype}
\usepackage{tikz}
\usepackage{tkz-euclide}
\usepackage{pgfplots}
\pgfplotsset{compat=1.18}
\usepackage{tabularx, booktabs, longtable}
\usepackage{enumitem}
\usepackage{hyperref}
\usepackage{parskip}
\usepackage{fancyhdr}
\usepackage{setspace}
\usepackage{xcolor}

\definecolor{accent}{HTML}{008080}
\definecolor{darkaccent}{HTML}{004D4D}
\definecolor{textgray}{HTML}{333333}
\definecolor{softgray}{HTML}{F3F5F5}
\definecolor{lightaccent}{HTML}{E6F2F2}

\providecommand{\tg}{\mathop{\mathrm{tg}}\nolimits}
\providecommand{\ctg}{\mathop{\mathrm{ctg}}\nolimits}
\providecommand{\arctg}{\mathop{\mathrm{arctg}}\nolimits}

\hypersetup{
    colorlinks=true,
    linkcolor=darkaccent,
    urlcolor=darkaccent
}

\onehalfspacing
\setlength{\parindent}{0pt}
\setlength{\parskip}{6pt}

\pagestyle{fancy}
\fancyhf{}
\fancyhead[L]{\small Чек-лист}
\fancyhead[R]{\small Тригонометрия и логарифмические неравенства}
\fancyfoot[C]{\small \thepage}
\renewcommand{\headrulewidth}{0.4pt}
\renewcommand{\headrule}{\hbox to\headwidth{\color{accent}\leaders\hrule height \headrulewidth\hfill}}

\setlist[itemize]{leftmargin=1.2cm,itemsep=2pt,topsep=4pt}
\setlist[enumerate]{leftmargin=1.2cm,itemsep=4pt,topsep=4pt}

\newcommand{\important}[1]{\textcolor{darkaccent}{\textbf{#1}}}
\newcommand{\ruleline}{\vspace{2mm}{\color{accent}\hrule height 0.5pt}\vspace{2mm}}

\newcommand{\signaxis}{
\begin{center}
\begin{tikzpicture}[scale=0.95]
    \draw[->, color=textgray] (-4.2,0) -- (4.8,0);
    \foreach \x/\lab in {-2/$-2$,0/$0$,3/$4$}
        \draw[color=textgray] (\x,0.08) -- (\x,-0.08) node[below=3pt] {\small \lab};
    \filldraw[fill=accent, draw=accent] (-2,0) circle (2pt);
    \filldraw[fill=white, draw=darkaccent, line width=1pt] (0,0) circle (2pt);
    \filldraw[fill=accent, draw=accent] (3,0) circle (2pt);
    \node[above] at (-3.1,0.2) {\small $-$};
    \node[above] at (-1.0,0.2) {\small $+$};
    \node[above] at (1.5,0.2) {\small $-$};
    \node[above] at (3.9,0.2) {\small $+$};
    \draw[line width=1.4pt, color=accent] (-4.0,0.35) -- (-2,0.35);
    \draw[line width=1.4pt, color=accent] (0,0.35) -- (3,0.35);
    \node[above, color=darkaccent] at (-3.0,0.45) {\small берём};
    \node[above, color=darkaccent] at (1.5,0.45) {\small берём};
\end{tikzpicture}
\end{center}
}

\begin{document}

\begin{titlepage}
\thispagestyle{empty}
\begin{center}
\vspace*{2.2cm}

{\Large \textcolor{accent}{\textbf{Чек-лист}}}

\vspace{8mm}

{\huge \textbf{Комплексное повторение ЕГЭ}}

\vspace{4mm}

{\Large \textbf{Тригонометрические уравнения и логарифмические неравенства}}

\vspace{18mm}

{\large Формулы, однородные уравнения, отбор корней, ОДЗ и метод интервалов}

\vspace{22mm}

{\large Лёвин Артём Александрович эксклюзивно для Николь Саркисьянц}

\vspace{8mm}

{\large \today}

\vfill

\begin{tikzpicture}[remember picture, overlay]
    \draw[line width=1.1pt, color=accent]
        ([xshift=2.3cm,yshift=2.2cm]current page.south west)
        .. controls ([xshift=6.0cm,yshift=3.3cm]current page.south west)
        and ([xshift=-6.0cm,yshift=1.3cm]current page.south east)
        .. ([xshift=-2.3cm,yshift=2.4cm]current page.south east);
    \draw[line width=0.5pt, color=darkaccent]
        ([xshift=3.0cm,yshift=1.55cm]current page.south west)
        .. controls ([xshift=7.4cm,yshift=2.25cm]current page.south west)
        and ([xshift=-7.0cm,yshift=2.7cm]current page.south east)
        .. ([xshift=-3.0cm,yshift=1.65cm]current page.south east);
\end{tikzpicture}

\vspace*{1.2cm}
\end{center}
\end{titlepage}

\section*{1. Главная идея занятия}

В комплексном повторении важно не просто знать формулы, а понимать, \important{зачем} они применяются.

\important{Главный принцип в тригонометрии:} если в выражении есть разные углы или квадраты, нужно искать формулу, которая приведёт выражение к более простому виду и, по возможности, к одинаковым углам.

\important{Главный принцип в неравенствах:} к уравнению, корням числителя и методу интервалов переходят только после того, как выражение приведено к нормальному виду:
\[
\frac{P(x)}{Q(x)}\le 0,
\qquad
\frac{P(x)}{Q(x)}\ge 0,
\qquad
\frac{P(x)}{Q(x)}<0,
\qquad
\frac{P(x)}{Q(x)}>0.
\]

\important{Отдельная привычка для самопроверки:} после каждого раскрытия скобок проверяйте знаки. Ошибка в одном минусе может полностью изменить ответ.

\section*{2. Формулы, которые чаще всего спасают тригонометрию}

\subsection*{Формулы понижения степени}

Если в задаче появились квадраты синуса или косинуса, почти всегда нужно вспомнить:
\[
\sin^2\alpha=\frac{1-\cos 2\alpha}{2},
\qquad
\cos^2\alpha=\frac{1+\cos 2\alpha}{2}.
\]

\important{Зачем это нужно?} Эти формулы убирают квадрат и переводят задачу к углу $2\alpha$.

\subsection*{Формулы разности}

\[
\cos(\alpha-\beta)=\cos\alpha\cos\beta+\sin\alpha\sin\beta,
\]
\[
\sin(\alpha-\beta)=\sin\alpha\cos\beta-\cos\alpha\sin\beta.
\]

\important{Полезный пример:}
\[
\sin^2\left(x-\frac{\pi}{4}\right)
=
\frac{1-\cos\left(2x-\frac{\pi}{2}\right)}{2}.
\]

Так как
\[
\cos\left(2x-\frac{\pi}{2}\right)
=
\cos 2x\cdot \cos\frac{\pi}{2}
+
\sin 2x\cdot \sin\frac{\pi}{2}
=
\sin 2x,
\]
то
\[
\boxed{
\sin^2\left(x-\frac{\pi}{4}\right)=\frac{1-\sin 2x}{2}.
}
\]

\subsection*{Формулы двойного угла}

\[
\cos 2x=\cos^2 x-\sin^2 x,
\]
\[
\sin 2x=2\sin x\cos x.
\]

\subsection*{Основное тригонометрическое тождество}

\[
\sin^2 x+\cos^2 x=1.
\]

\important{Важная идея.} Иногда единицу полезно заменить на тригонометрическое выражение:
\[
1=\sin^2 x+\cos^2 x.
\]

Это помогает получить однородное уравнение.

\section*{3. Образцовый тригонометрический пример}

Рассмотрим уравнение:
\[
\cos^2x+\sin^2\left(x-\frac{\pi}{4}\right)=\frac12.
\]

По формулам понижения степени:
\[
\cos^2x=\frac{1+\cos2x}{2},
\]
\[
\sin^2\left(x-\frac{\pi}{4}\right)=\frac{1-\sin2x}{2}.
\]

Подставляем:
\[
\frac{1+\cos2x}{2}
+
\frac{1-\sin2x}{2}
=
\frac12.
\]

Домножаем на $2$:
\[
1+\cos2x+1-\sin2x=1.
\]

Получаем:
\[
\cos2x-\sin2x=-1.
\]

Теперь раскрываем через формулы двойного угла:
\[
\cos^2x-\sin^2x-2\sin x\cos x
=
-\left(\sin^2x+\cos^2x\right).
\]

Переносим всё влево:
\[
\cos^2x-\sin^2x-2\sin x\cos x+\sin^2x+\cos^2x=0.
\]

Приводим подобные:
\[
2\cos^2x-2\sin x\cos x=0.
\]

Выносим общий множитель:
\[
2\cos x(\cos x-\sin x)=0.
\]

Получаем совокупность:
\[
\left[
\begin{array}{l}
\cos x=0,\\
\cos x-\sin x=0.
\end{array}
\right.
\]

Отсюда:
\[
x=\frac{\pi}{2}+\pi k,
\qquad k\in\mathbb Z,
\]
или
\[
\tg x=1
\quad\Longleftrightarrow\quad
x=\frac{\pi}{4}+\pi k,
\qquad k\in\mathbb Z.
\]

Ответ:
\[
\boxed{
x=\frac{\pi}{2}+\pi k
\quad\text{или}\quad
x=\frac{\pi}{4}+\pi k,
\qquad k\in\mathbb Z.
}
\]

\section*{4. Однородные тригонометрические уравнения}

\subsection*{Однородное уравнение первого порядка}

Типичный вид:
\[
a\sin x+b\cos x=0.
\]

Пример:
\[
\cos x-\sin x=0.
\]

Если $\cos x=0$, то равенство не выполняется. Значит, можно делить на $\cos x$:
\[
\frac{\cos x}{\cos x}-\frac{\sin x}{\cos x}=0.
\]

Получаем:
\[
1-\tg x=0,
\]
\[
\tg x=1,
\]
\[
x=\frac{\pi}{4}+\pi k,\qquad k\in\mathbb Z.
\]

\important{Почему важно проверить деление?} Делить на выражение можно только тогда, когда оно не равно нулю на рассматриваемой ветви решения.

\subsection*{Однородное уравнение второго порядка}

Типичный вид:
\[
a\sin^2x+b\sin x\cos x+c\cos^2x=0.
\]

Пример:
\[
\sin^2x-3\sin x\cos x+2\cos^2x=0.
\]

Если $\cos x=0$, то левая часть равна $\sin^2x=1$, значит, это не корень. Поэтому можно делить на $\cos^2x$:
\[
\tg^2x-3\tg x+2=0.
\]

Дальше решается обычное квадратное уравнение относительно $\tg x$.

\important{Неоднородное уравнение} нельзя решать таким же делением. Например:
\[
3+\sin x+5\cos x=0
\]
не является однородным, потому что слева есть отдельное число $3$.

\section*{5. Отбор корней на промежутке}

Если в условии дан промежуток, например
\[
x\in[3\pi;4\pi],
\]
то после нахождения общих корней нужно отобрать подходящие.

\subsection*{Пример отбора}

Пусть
\[
x=\frac{\pi}{2}+\pi k,
\qquad k\in\mathbb Z,
\]
и нужно отобрать корни на промежутке:
\[
3\pi\le x\le 4\pi.
\]

Подставляем формулу корня:
\[
3\pi\le \frac{\pi}{2}+\pi k\le 4\pi.
\]

Домножаем на $2$:
\[
6\pi\le \pi+2\pi k\le 8\pi.
\]

Вычитаем $\pi$:
\[
5\pi\le 2\pi k\le 7\pi.
\]

Делим на $2\pi$:
\[
\frac52\le k\le \frac72.
\]

Так как $k\in\mathbb Z$, подходит только:
\[
k=3.
\]

Тогда:
\[
x=\frac{\pi}{2}+3\pi=\frac{7\pi}{2}.
\]

\important{Практическое правило.} Отбор через двойное неравенство медленнее, чем по окружности, но он очень надёжен и хорошо защищает от ошибок невнимательности.

\section*{6. Логарифмические неравенства: ОДЗ, замена и метод интервалов}

\subsection*{Обязательные ограничения}

Для логарифма:
\[
\log_a f(x)
\]
должны выполняться условия:
\[
a>0,\qquad a\ne1,\qquad f(x)>0.
\]

Если логарифм стоит в знаменателе, дополнительно:
\[
\log_a f(x)\ne0.
\]

\subsection*{Ключевые свойства логарифмов}

\[
\log_a\frac{u}{v}=\log_a u-\log_a v,
\qquad u>0,\quad v>0.
\]

\[
\log_a u^m=m\log_a|u|.
\]

\important{Осторожно с модулем:}
\[
\log_a x^2=2\log_a|x|.
\]

Если по ОДЗ уже известно, что $x>0$, то:
\[
|x|=x,
\]
и поэтому:
\[
\log_a x^2=2\log_a x.
\]

\subsection*{Образцовый логарифмический пример}

Рассмотрим неравенство:
\[
\frac{9}{\log_3x}-\log_3\frac{9}{x}
\le
\frac{34}{\log_3x^2}.
\]

ОДЗ:
\[
x>0,
\qquad
\log_3x\ne0.
\]

То есть:
\[
x>0,
\qquad
x\ne1.
\]

Преобразуем логарифмы:
\[
\log_3\frac{9}{x}
=
\log_39-\log_3x
=
2-\log_3x.
\]

Так как $x>0$:
\[
\log_3x^2=2\log_3x.
\]

Подставляем:
\[
\frac{9}{\log_3x}
-
\left(2-\log_3x\right)
\le
\frac{34}{2\log_3x}.
\]

Сокращаем справа:
\[
\frac{9}{\log_3x}
-
\left(2-\log_3x\right)
\le
\frac{17}{\log_3x}.
\]

Вводим замену:
\[
t=\log_3x,
\qquad t\ne0.
\]

Получаем:
\[
\frac9t-(2-t)-\frac{17}{t}\le0.
\]

Приводим к общему знаменателю:
\[
\frac{9-2t+t^2-17}{t}\le0.
\]

То есть:
\[
\frac{t^2-2t-8}{t}\le0.
\]

Разложим числитель:
\[
t^2-2t-8=(t-4)(t+2).
\]

Получаем рациональное неравенство:
\[
\frac{(t-4)(t+2)}{t}\le0.
\]

Критические точки:
\[
t=-2,\qquad t=0,\qquad t=4.
\]

Точка $t=0$ выколота, потому что стоит в знаменателе.

\signaxis

Отсюда:
\[
t\in(-\infty;-2]\cup(0;4].
\]

Возвращаемся к замене:
\[
\log_3x\le -2
\quad\text{или}\quad
0<\log_3x\le4.
\]

Так как основание $3>1$, знак неравенства не меняется:
\[
0<x\le\frac19
\quad\text{или}\quad
1<x\le81.
\]

Ответ:
\[
\boxed{
x\in\left(0;\frac19\right]\cup(1;81].
}
\]

\section*{7. Типовые ошибки}

\begin{enumerate}
    \item \important{Переходить к уравнению слишком рано.}

    Сначала нужно привести неравенство к нормальному виду, раскрыть скобки, привести подобные и общий знаменатель. Только потом можно искать нули числителя и знаменателя.

    \item \important{Домножать неравенство на переменную без учёта знака.}

    Нельзя просто домножить неравенство на $t$, если неизвестно, положительно оно или отрицательно. Надёжный путь:
    \[
    \text{привести к общему знаменателю и решить методом интервалов.}
    \]

    \item \important{Забывать модуль в формуле $\log_a x^2$.}

    Корректно:
    \[
    \log_a x^2=2\log_a|x|.
    \]

    Без модуля можно писать только после анализа ОДЗ.

    \item \important{Забывать запрещённые значения.}

    Если переменная стоит в знаменателе, соответствующая точка всегда выколота.

    \item \important{Терять знак при раскрытии скобок.}

    Особенно опасно:
    \[
    -\left(2-\log_3x\right)=-2+\log_3x.
    \]

    \item \important{Не проверять деление в однородных уравнениях.}

    Перед делением на $\cos x$ или $\cos^2x$ нужно понимать, не теряются ли корни.
\end{enumerate}

\section*{8. Мини-алгоритмы}

\subsection*{Тригонометрическое уравнение}

\begin{enumerate}
    \item Посмотрите, что мешает: квадраты, разные углы, сумма синуса и косинуса.
    \item Подберите формулу: понижение степени, разность, двойной угол.
    \item Добейтесь одинаковых углов или более простого выражения.
    \item Если появилось произведение, переходите к совокупности.
    \item Если уравнение однородное, проверьте возможность деления.
    \item Найдите общие корни.
    \item Если дан промежуток, выполните отбор через двойное неравенство.
\end{enumerate}

\subsection*{Логарифмическое неравенство}

\begin{enumerate}
    \item Выпишите ОДЗ.
    \item Разложите логарифмы по свойствам.
    \item Если возможно, введите замену.
    \item Не домножайте на переменную неизвестного знака.
    \item Приведите выражение к рациональному неравенству.
    \item Решите методом интервалов.
    \item Вернитесь к исходной переменной.
    \item Пересеките ответ с ОДЗ.
\end{enumerate}

\newpage

\section*{Тренировочный блок}

\textbf{Инструкция.} Решите задания. В тригонометрии отдельно отмечайте, какую формулу Вы применяете. В логарифмических неравенствах обязательно выписывайте ОДЗ.

\subsection*{Средний уровень}

\begin{enumerate}
    \item Упростите выражение:
    \[
    \sin^2\left(x-\frac{\pi}{4}\right).
    \]

    \item Решите уравнение:
    \[
    2\cos x(\cos x-\sin x)=0.
    \]

    \item Решите рациональное неравенство:
    \[
    \frac{(t-4)(t+2)}{t}\le0.
    \]
\end{enumerate}

\subsection*{Выше среднего уровня}

\begin{enumerate}[resume]
    \item Решите тригонометрическое уравнение:
    \[
    \cos^2x+\sin^2\left(x-\frac{\pi}{4}\right)=\frac12.
    \]

    \item Отберите корни формулы
    \[
    x=\frac{\pi}{4}+\pi k,\qquad k\in\mathbb Z,
    \]
    на промежутке:
    \[
    x\in\left[\frac{13\pi}{2};\frac{15\pi}{2}\right].
    \]

    \item Решите уравнение при условии $x>0$:
    \[
    \log_3x^2=4.
    \]

    \item Решите логарифмическое неравенство:
    \[
    \frac{9}{\log_3x}-\log_3\frac{9}{x}
    \le
    \frac{34}{\log_3x^2}.
    \]

    \item Решите однородное тригонометрическое уравнение:
    \[
    \sin^2x-3\sin x\cos x+2\cos^2x=0.
    \]

    \item Решите логарифмическое неравенство:
    \[
    \log_3|x-2|\le2.
    \]
\end{enumerate}

\subsection*{Сложный уровень}

\begin{enumerate}[resume]
    \item Решите уравнение и выполните отбор корней на указанном промежутке:
    \[
    \cos^2x+\sin^2\left(x-\frac{\pi}{4}\right)=\frac12,
    \qquad
    x\in\left[\frac{13\pi}{2};\frac{15\pi}{2}\right].
    \]
\end{enumerate}

\newpage

\section*{Ответы для самопроверки}

\begin{center}
\renewcommand{\arraystretch}{1.45}
\begin{tabularx}{\linewidth}{>{\bfseries}p{0.12\linewidth}X}
\toprule
№ & Ответ \\
\midrule
1 & $\displaystyle \frac{1-\sin2x}{2}$ \\
2 & $\displaystyle x=\frac{\pi}{2}+\pi k$ или $\displaystyle x=\frac{\pi}{4}+\pi k,\quad k\in\mathbb Z$ \\
3 & $\displaystyle t\in(-\infty;-2]\cup(0;4]$ \\
4 & $\displaystyle x=\frac{\pi}{2}+\pi k$ или $\displaystyle x=\frac{\pi}{4}+\pi k,\quad k\in\mathbb Z$ \\
5 & $\displaystyle k=7,\quad x=\frac{29\pi}{4}$ \\
6 & $\displaystyle x=9$ \\
7 & $\displaystyle x\in\left(0;\frac19\right]\cup(1;81]$ \\
8 & $\displaystyle x=\frac{\pi}{4}+\pi k$ или $\displaystyle x=\arctg2+\pi k,\quad k\in\mathbb Z$ \\
9 & $\displaystyle x\in[-7;11]\setminus\{2\}$ \\
10 & $\displaystyle x\in\left\{\frac{13\pi}{2};\frac{29\pi}{4};\frac{15\pi}{2}\right\}$ \\
\bottomrule
\end{tabularx}
\end{center}

\end{document}